% !PRINTED 457
\weberpage{457}
Bezeichnen wir die Coëfficienten von $\varphi(x)$, die, wie wir gesehen haben, Grössen in $\Omega$ sind, mit $c_0, c_1 \ldots c_{n-1}$, so kann also jede Grösse $\Theta$ in $\Omega(\alpha)$ in der Form dargestellt werden:

\[
\Theta = c_0 + c_1\alpha + c_2\alpha^2 + \cdots + c_{n-1}\alpha^{n-1}.
\]

Diese Darstellung ist nur auf \gesperrt{eine Art} möglich; denn ist gleichzeitig

\weq{6}{\Theta = c_0' + c_1'\alpha + c_2'\alpha^2 + \cdots + c_{n-1}'\alpha^{n-1},}

so verschwindet die Function $(n-1)^{\text{ten}}$ Grades

\[
c_0 - c_0' + (c_1 - c_1')x + (c_2 - c_2')x^2 + \cdots + (c_{n-1} - c_{n-1}')x^{n-1}
\]

für $x = \alpha$. Das ist aber nach dem Satz III des vorigen Paragraphen nur möglich, wenn

\[
c_0 = c_0', \quad c_1 = c_1', \ldots c_{n-1} = c_{n-1}'
\]

ist.

Wenn wir die Grössen des Körpers $\Omega$ als die gegebenen, seine Grössen als rational, den Körper $\Omega(\alpha)$ als den daraus abgeleiteten betrachten, so bezeichnen wir eine rationale Function in $\Omega$ auch kurz als eine \gesperrt{rationale Function}, und können demnach auch sagen:

\gesperrt{Der Körper $\Omega(\alpha)$ besteht aus allen rationalen Functionen von $\alpha$.}

\begin{center}
\S.\,143.

\gesperrt{Gleichzeitige Adjunction mehrerer algebraischer Grössen.}
\end{center}

Einer genaueren Untersuchung des Körpers $\Omega(\alpha)$ schicken wir einen allgemeinen, sehr einfachen, aber doch folgereichen Satz voraus.

1. \gesperrt{Sind $\Phi_1(x,y,z\ldots)$, $\Phi_2(x,y,z\ldots)$, $\Phi_3(x,y,z\ldots)\ldots$ ganze rationale Functionen der Veränderlichen $x,y,z\ldots$ mit beliebigen Coëfficienten, die in keiner der Functionen alle zugleich verschwinden, so kann man für die Veränderlichen auf unendlich viele Arten solche rationale Zahlwerthe setzen, dass keine der Functionen $\Phi_1, \Phi_2, \Phi_3\ldots$ verschwindet.}
